% Generated draft: requires mathematical and source review.
\documentclass[11pt]{article}
\usepackage[margin=25mm]{geometry}
\usepackage{fontspec}
\setmainfont{Latin Modern Roman}
\usepackage{amsmath,amssymb,mathtools,unicode-math}
\setmathfont{Latin Modern Math}
\usepackage{enumitem,needspace,xurl,hyperref}
\setlength{\parindent}{0pt}
\setlength{\parskip}{5pt}
\setlength{\emergencystretch}{3em}
\newcommand{\N}{\mathbb{N}}
\newcommand{\Z}{\mathbb{Z}}
\newcommand{\Q}{\mathbb{Q}}
\newcommand{\R}{\mathbb{R}}
\newcommand{\C}{\mathbb{C}}
\newcommand{\F}{\mathbb{F}}
\newcommand{\A}{\mathbb{A}}
\newcommand{\PP}{\mathbb{P}}
\newcommand{\E}{\mathbb{E}}
\newcommand{\eps}{\varepsilon}
\newcommand{\dd}{\,\mathrm{d}}
\begin{document}
\section*{UnsolvedMath ID: 11000271}
\textbf{Source problem code:} catalogue entry
\subsection{Definitions and mathematical statement}
Let \(G=\operatorname{Mod}(S_g)\) be the mapping class group of a closed surface in the source's convention. Its virtual cohomological dimension is the cohomological dimension of a torsion-free finite-index subgroup; in the usual closed orientable genus-\(g\geq2\) case it is \(4g-5\). The question asks for a cocompact classifying space denoted ``\(EG\)'' whose dimension equals this virtual dimension.

The bar/underline convention is missing from the extraction. The natural proper-action interpretation seeks a model \(X\) for \(\underline EG=E_{\mathrm{Fin}}G\), a \(G\)-CW complex such that
\[X^H\text{ is contractible for every finite }H\leq G,\qquad
X^H=\varnothing\text{ for every infinite }H\leq G,\]
with \(X/G\) compact and \(\dim X=\operatorname{vcd}(G)\). Thus, in the indicated \(g\geq2\) setting, the target dimension is \(4g-5\). This combines the minimum-dimensional target with cocompactness, rather than asking for either one alone.

If the notation instead means an ordinary free contractible \(EG\), that is a different request: a group with torsion has no finite-dimensional free contractible model, whereas mapping class groups generally have torsion. A torsion-free finite-index subgroup's classifying space is another distinct object. The original notation must therefore be checked before assigning a precise target. Likewise \(E_{\mathrm{VCyc}}G\), for the virtually cyclic family mentioned vaguely in the summary, has different fixed-point requirements and cannot be substituted for \(E_{\mathrm{Fin}}G\). The surface genus/extended-group conventions and intended uniformity of construction are not stated.
\subsection{Short English statement}
Construct a cocompact proper-action classifying space for the closed-surface mapping class group in virtual cohomological dimension, if that is the intended \(EG\) notation. Ordinary free \(EG\) and the virtually-cyclic-family version are different targets.
\subsection{Sources}
\begin{itemize}
\item Local UnsolvedMath registry, permanent ID \(11000271\), exact source code \texttt{catalogue entry}.
\item Farb, \emph{Problems on Mapping Class Groups and Related Topics}, Question 2.3, PDF page 330: \url{https://www.math.uchicago.edu/~farb/papers/mcgbook.pdf}. Registry access 2026-07-29; pdf-font-structure extraction.
\item The summary gives only an incomplete ``Bianchi--...'' lead and a virtually-cyclic-classifying-space analogy, without a usable citation or exact resolving theorem.
\end{itemize}
\end{document}
